\chapter{Hashimoto Thyroiditis}

\section*{Definition and Overview}
Hashimoto thyroiditis, also called Hashimoto's disease or chronic autoimmune thyroiditis, is an autoimmune condition in which the immune system attacks the thyroid gland, gradually destroying its ability to produce thyroid hormone. It is the most common cause of hypothyroidism (an underactive thyroid) in Australia. The condition is more common in women and often develops slowly over years, with the gland sometimes enlarging into a goitre before shrinking as it is destroyed. Most people develop hypothyroidism and need lifelong thyroid hormone replacement, which is safe and effective. Because the condition is so common, thyroid blood tests are routine in Australia.

\section*{Epidemiology}
Hashimoto thyroiditis is the most common cause of hypothyroidism, which affects around 3--7% of the population. It is five to ten times more common in women than men and typically begins between ages 30 and 50, though it can occur at any age, including in children. The prevalence increases with age, affecting up to 10--15% of older women. It is more common in people with a family history of thyroid or other autoimmune disease, and it often coexists with other autoimmune conditions such as coeliac disease, type 1 diabetes, and pernicious anaemia.

\section*{Aetiology and Pathophysiology}
Hashimoto thyroiditis is caused by an autoimmune attack on the thyroid gland.

\begin{itemize}
    \item \textbf{Autoantibodies:} the immune system produces antibodies against thyroid proteins, including thyroid peroxidase (TPO) and thyroglobulin
    \item \textbf{Chronic inflammation:} antibodies and immune cells infiltrate the gland, causing inflammation and gradual destruction of thyroid tissue
    \item \textbf{Hypothyroidism:} as the gland is destroyed, it can no longer produce enough thyroid hormone
    \item \textbf{Goitre:} early in the disease, the gland may enlarge as it tries to compensate
    \item \textbf{Genetic susceptibility:} certain HLA genes increase risk
    \item \textbf{Triggers:} viral infections, stress, pregnancy, and excess iodine can trigger onset in susceptible people
    \item \textbf{Other autoimmune disease:} affected people are prone to other autoimmune conditions
\end{itemize}

\section*{Clinical Features}
Symptoms of hypothyroidism develop gradually and are often non-specific.

\begin{itemize}
    \item Fatigue, weakness, and lethargy
    \item Weight gain despite normal eating
    \item Cold intolerance and feeling cold
    \item Constipation
    \item Dry skin and hair, and hair loss
    \item Depression, poor concentration, and memory problems
    \item Muscle aches, joint pain, and stiffness
    \item Heavy or irregular periods
    \item A goitre (enlarged thyroid), which can cause a feeling of fullness in the throat
    \item Slow heart rate and hoarse voice in advanced cases
    \item Many people have few or no symptoms initially
\end{itemize}

\section*{Differential Diagnosis}
Other causes of hypothyroidism and similar symptoms should be considered.

\begin{itemize}
    \item Other causes of hypothyroidism, including iodine deficiency, thyroid surgery, and radiation treatment
    \item Postpartum thyroiditis and other thyroid conditions
    \item Graves disease, which causes hyperthyroidism and can coexist
    \item Chronic fatigue syndrome and fibromyalgia
    \item Depression and anxiety
    \item Anaemia, including iron and B12 deficiency
    \item Other causes of goitre, including nodules and cancer
    \item Menopause and other hormonal changes
\end{itemize}

\section*{When to Seek Medical Care}
See a doctor if you have persistent symptoms of hypothyroidism, especially fatigue, weight gain, cold intolerance, or changes in periods. Thyroid tests are simple and part of routine health care.

\textbf{Call 000 (or local emergency services) for:}
\begin{itemize}
    \item Symptoms of myxoedema coma, a rare emergency: severe confusion, low body temperature, and slowed breathing in a person with severe hypothyroidism
    \item Chest pain or severe breathlessness
    \item A rapidly growing neck lump with difficulty breathing or swallowing
\end{itemize}

\section*{Diagnosis and Investigations}
Hashimoto thyroiditis is diagnosed with blood tests.

\begin{itemize}
    \item \textbf{Thyroid function tests:} elevated TSH with low free T4 confirms hypothyroidism
    \item \textbf{Thyroid antibodies:} elevated TPO antibodies confirm the autoimmune cause
    \item \textbf{Ultrasound:} may show a characteristic appearance of the gland and assess goitre and nodules
    \item \textbf{Thyroid function monitoring:} regular tests track the condition
    \item \textbf{Other blood tests:} for coexisting conditions, including vitamin B12 and iron
    \item \textbf{Antenatal screening:} thyroid function is checked in pregnancy for women at risk
\end{itemize}

\section*{Management}
Treatment is simple, safe, and lifelong for most people.

\begin{itemize}
    \item \textbf{Levothyroxine:} synthetic thyroid hormone replacement, taken daily, usually for life
    \item \textbf{Dose adjustment:} doses are adjusted based on regular blood tests
    \item \textbf{Monitoring:} thyroid function is checked initially, then at intervals and after dose changes
    \item \textbf{Timing of the medicine:} levothyroxine is best taken on an empty stomach, typically before breakfast
    \item \textbf{Subclinical hypothyroidism:} mildly raised TSH with normal free T4 may or may not be treated, depending on the situation
    \item \textbf{Goitre management:} hormone replacement may shrink an early goitre
    \item \textbf{Pregnancy:} thyroid dose usually increases in pregnancy, and monitoring is intensified
    \item \textbf{Lifestyle:} a healthy diet with adequate iodine and selenium
\end{itemize}

\section*{Complications}
\begin{itemize}
    \item Untreated hypothyroidism: heart disease, high cholesterol, and, in pregnancy, effects on the baby's development
    \item Myxoedema coma, a rare life-threatening emergency
    \item Goitre causing pressure symptoms
    \item Increased risk of other autoimmune diseases
    \item Slightly increased risk of thyroid lymphoma, which is very rare
    \item Effects of undertreatment or overtreatment with levothyroxine
    \item Psychological effects of chronic fatigue and other symptoms before diagnosis
\end{itemize}

\section*{Prognosis and Outcomes}
With levothyroxine treatment, people with Hashimoto thyroiditis lead normal, healthy lives. Symptoms resolve within weeks of treatment starting, and blood tests confirm adequate dosing. Treatment is lifelong, but it is simple and safe when monitored. Untreated hypothyroidism causes progressive symptoms and health effects, so diagnosis and treatment matter. The prognosis is excellent for people who take their medicine consistently and attend monitoring.

\section*{Prevention and Self-Care}
\begin{itemize}
    \item Take levothyroxine at the same time each day, on an empty stomach
    \item Do not stop treatment without medical advice
    \item Attend regular blood tests and dose reviews
    \item Avoid taking calcium, iron, or antacids at the same time as levothyroxine
    \item Eat a balanced diet with adequate iodine, but do not take iodine supplements unless advised
    \item Report new or worsening symptoms promptly
    \item If pregnant or planning pregnancy, ensure thyroid function is checked and managed
    \item Have blood tests for other autoimmune conditions if symptoms suggest them
    \item Carry a list of your medicines and thyroid status
\end{itemize}

\section*{Sources and Further Reading}
The following reputable sources provide further information for healthcare students and patients seeking reliable, current advice about Hashimoto thyroiditis.

\begin{itemize}
    \item Australian Thyroid Foundation. \textit{Hashimoto thyroiditis information.} https://www.thyroidfoundation.org.au/
    \item Healthdirect Australia. \textit{Hypothyroidism and Hashimoto disease.} https://www.healthdirect.gov.au/
    \item Jean Hailes for Women's Health. \textit{Thyroid conditions.} https://www.jeanhailes.org.au/
    \item NPS MedicineWise. \textit{Levothyroxine: using it well.} https://www.nps.org.au/
    \item American Thyroid Association. \textit{Hashimoto's thyroiditis.} https://www.thyroid.org/
    \item NHS. \textit{Underactive thyroid (hypothyroidism).} https://www.nhs.uk/
\end{itemize}
